%% ma: L11-B27-0004
%% lop: 11
%% chuong: 7
%% bai: 27
%% dang: 11.27.4
%% loai: TLN
%% muc: VD
%% dap_an: "5"
%% nguon: "(NBV)-ĐỀ SỐ 3-MỖI NGÀY 1 ĐỀ THI 2025.pdf (D03, MỖI NGÀY 1 ĐỀ - Đề số 3), Phần III Câu 2"
%% kien_thuc: "Bài 25 (hai mặt phẳng vuông góc), Bài 24 (góc giữa đường thẳng và mặt phẳng), Bài 27 (thể tích khối chóp)"
%% trang_thai: chinh-thuc
%% ghi_chu: "Đáp án gốc 5 đúng. Bỏ hình (đề không có hình)"
\begin{ex}
Cho hình chóp $S.ABCD$ có đáy $ABCD$ là hình vuông cạnh bằng $2a$. Tam giác $SAB$ cân tại $S$ và $(SAB)\perp(ABCD)$. Biết thể tích của khối chóp $S.ABCD$ là $\dfrac{4a^3}{3}$. Gọi $\alpha$ là góc giữa $SC$ và $(ABCD)$. Khi đó $\tan\alpha=\dfrac{\sqrt n}{5}$, $n\in\mathbb N$. Giá trị của $n$ là bao nhiêu?
\shortans{5}
\loigiai{Gọi $H$ là trung điểm $AB$. Tam giác $SAB$ cân tại $S$ nên $SH\perp AB$; mà $(SAB)\perp(ABCD)$ nên $SH\perp(ABCD)$.
$V=\dfrac13\cdot(2a)^2\cdot SH=\dfrac{4a^3}{3}\Rightarrow SH=a$.
Góc giữa $SC$ và $(ABCD)$ là $\alpha=\widehat{SCH}$, với $HC=\sqrt{HB^2+BC^2}=\sqrt{a^2+4a^2}=a\sqrt5$.
$\tan\alpha=\dfrac{SH}{HC}=\dfrac{1}{\sqrt5}=\dfrac{\sqrt5}{5}$, suy ra $n=5$.}
\end{ex}
